\section{Актуальность, цель и задачи исследования}

TODO2

Предметом работы является разработка программного модуля «НейроКОМПАС-3D», включающего обучаемые модели, средства подготовки данных, исполнительные сценарии инференса и интеграционный контур взаимодействия с КОМПАС-3D. В задаче присутствуют три взаимосвязанные подсистемы: дата-инженерный контур; нейросетевой контур; контур интеграции с CAD-средой.

Целью является разработка программного модуля «НейроКОМПАС-3D», включающего обучаемые модели, средства подготовки данных, исполнительные сценарии инференса и интеграционный контур взаимодействия с КОМПАС-3D. В задаче присутствуют три взаимосвязанные подсистемы: дата-инженерный контур; нейросетевой контур; контур интеграции с CAD-средой.

Задачи:

Определены поддерживаемые операции, формат набора данных, предложены CNN/Transformer/GNN-подходы. Разработаны требования к набору данных для машинного обучения. Сформированы первичные наборы данных для обучения нейронной сети. Разработана архитектура нейронной сети, пригодная для решения поставленных задач.

Разработан dataloader, выполнен препроцессинг, обучена базовая модель, подтверждена концептуальная возможность. Проведено предварительное обучение на основании подготовленного набора данных. Сформирован отчет с техническими показателями обучения нейронной сети и предложениями по дальнейшему улучшению показателей, в том числе с помощью корректировки архитектуры модели и исходных данных.

Подготовлены модели для Steps.txt и геометрии, оформлен комплект ПО и весов. Разработаны нейронные сети, которые способны формировать один шаг построения на основе двух моделей в определённом на этапе 1 формате (начальной и конечной), отличающихся одним элементарным построением  из набора элементарных операций, определённом на этапе 1.

Обоснован и принят сегментационный подход, обработан подготовленный специалистами Заказчика новый модифицированный крупный датасет, доработаны и обучены нейросети классификаторов и сегментаторов для каждого типа выбранных операций, разработан модуль main.exe (интегрируемый с КОМПАС-3D). Разработаны нейронные сети, которые способны пошагово (в итерационном взаимодействии с КОМПАС-3D) формировать промежуточные ветви дерева построения на основе двух моделей в определённом на этапе 1 формате (начальной и конечной), отличающихся несколькими элементарными построениями из набора элементарных операций, определённом на этапе 1.
